%!TEX root = ../exa-ma-d7.3.tex
\section{Introduction}
\label{sec:intro}
%---------------------------------------

\subsection{Context and Objectives}
The Exa-MA project (\textit{Méthodes et Algorithmes à l'Exascale}) is a key component of the national PEPR NumPEx program, targeting breakthrough research and software development to enable exascale-level computations in multiple scientific and industrial domains. 
Exa-MA focuses on advancing numerical methods and algorithms that address core exascale bottlenecks (B1--B13), including communication avoidance, handling of heterogeneous hardware, resiliency, and data-intensive workflows.

Within Exa-MA, \textbf{Work Package 7 (WP7)} is dedicated to ensuring that novel exascale solutions reach a wide audience of practitioners and stakeholders, and \textbf{Task 7.4 (Training)} specifically deals with creating, consolidating, and disseminating training materials. 
These materials are designed to support both new and experienced users who wish to integrate Exa-MA's cutting-edge methodologies and software stacks into their research or industrial workflows.

\textbf{Deliverable D7.3}, as part of WP7, provides an \textit{initial} snapshot of the project's training-related outputs and plans. It outlines:
\begin{itemize}
    \item The overall \textit{training strategy} (Section~\ref{sec:strategy}) to equip various audiences with exascale computing skills.
    \item The \textit{current status} of training materials (Section~\ref{sec:current}), including existing tutorials, workshops, and potential academic collaborations.
    \item Future steps to \textit{expand and refine} these materials, thereby aligning with the evolving needs of Exa-MA (Section~\ref{sec:next-steps-software}).
\end{itemize}

This document will be updated at subsequent milestones (M24, M36, \dots) to capture ongoing developments, feedback, and best practices arising from Exa-MA research and deployment activities.
